% Generated draft: requires mathematical and source review.
\documentclass[11pt]{article}
\usepackage[margin=25mm]{geometry}
\usepackage{fontspec}
\setmainfont{Latin Modern Roman}
\usepackage{amsmath,amssymb,mathtools,unicode-math}
\setmathfont{Latin Modern Math}
\usepackage{enumitem,needspace,xurl,hyperref}
\setlength{\parindent}{0pt}
\setlength{\parskip}{5pt}
\setlength{\emergencystretch}{3em}
\newcommand{\N}{\mathbb{N}}
\newcommand{\Z}{\mathbb{Z}}
\newcommand{\Q}{\mathbb{Q}}
\newcommand{\R}{\mathbb{R}}
\newcommand{\C}{\mathbb{C}}
\newcommand{\F}{\mathbb{F}}
\newcommand{\A}{\mathbb{A}}
\newcommand{\PP}{\mathbb{P}}
\newcommand{\E}{\mathbb{E}}
\newcommand{\eps}{\varepsilon}
\newcommand{\dd}{\,\mathrm{d}}
\begin{document}
\section*{UnsolvedMath ID: 700622}
\textbf{Source problem code:} catalogue entry
\subsection{Definitions and mathematical statement}
For a group \(G\), let \(F(X)\) be the free group on the countably infinite set \(X=\{x_1,x_2,\ldots\}\), and form the free product \(G*F(X)\). A word \(w\in G*F(X)\) is a mixed identity of \(G\) if every homomorphism \(\varphi:G*F(X)\to G\) whose restriction to \(G\) is the identity satisfies \(\varphi(w)=1_G\). Thus variables may be replaced independently by arbitrary elements of \(G\), while coefficients from \(G\) retain their specified values. Write
\[
\operatorname{MI}(G)=\bigcap_{\substack{\varphi:G*F(X)\to G\\ \varphi|_G=\mathrm{id}_G}}\ker\varphi.
\]
Use the following equational interpretation of a finite basis: a finite set \(B\subseteq\operatorname{MI}(G)\) is a basis if
\[
\operatorname{MI}(G)=\left\langle\!\left\langle
\alpha(b)\;\middle|\; b\in B,\ \alpha\in\operatorname{End}(G*F(X)),\ \alpha|_G=\mathrm{id}_G
\right\rangle\!\right\rangle_{G*F(X)},
\]
where the double brackets denote normal closure. This allows substitution of words with coefficients for variables, together with the usual consequences of group equations.

\textbf{Problem (V. S. Anashin, 1986).} Construct a finite group \(G\) for which \(\operatorname{MI}(G)\) has no finite basis. In addition, determine whether the group referred to in R. Bryant, \emph{Bulletin of the London Mathematical Society} \textbf{14}, no.~2 (1982), 119--123, is such an example.

\textbf{Unresolved source qualifications.} The supplied record does not define mixed identities or the consequence relation used for a basis; the definitions above make explicit a conventional interpretation that requires confirmation against the original source. The record does not give Bryant's group construction or a presentation, so the supplementary question identifies its target only by the supplied bibliographic reference.
\subsection{Short English statement}
Find a finite group whose equations with coefficients from that group cannot all be deduced from finitely many such equations, using the interpretation specified above. Also determine whether the group constructed in the cited 1982 paper by Bryant provides an example. The supplied registry labels the problem open; this status has not been independently verified.
\subsection{Sources}
\begin{itemize}
\item Local UnsolvedMath registry, permanent ID \(700622\), exact source problem code \texttt{catalogue entry}, title ``Kourovka Notebook Problem 10.5.'' Supplied record updated \texttt{2026-10-16T14:00:00Z}; upstream revision \texttt{d371ee917ac9647fbb06cd6fc529deec6f3dda82}. The registry status is supplied as \texttt{open}.
\item \emph{The Kourovka Notebook (Unsolved Problems in Group Theory)}, problem 10.5, proposed by V. S. Anashin, 10th issue (1986). The supplied background also identifies the 21st issue edition (2026), arXiv:\texttt{1401.0300v48}. The collection description names E. I. Khukhro and V. D. Mazurov as editors. This attribution and edition information are supplied source data, not independently verified here.
\item R. Bryant, \emph{Bulletin of the London Mathematical Society} \textbf{14}, no.~2 (1982), 119--123. This is the reference supplied for the group in the supplementary question; no article title or construction is included in the record.
\end{itemize}
\end{document}
